\documentclass[11pt]{article}
\usepackage{url}
\usepackage{alltt}
\usepackage{bm}
\linespread{1}
\textwidth 6.5in
\oddsidemargin 0.in
\addtolength{\topmargin}{-1in}
\addtolength{\textheight}{2in}

\usepackage{amsmath}
\usepackage{amssymb}
\usepackage{hyperref}

\begin{document}


\begin{center}
\Large
STA 214 Homework 3\\
\normalsize
\vspace{5mm}
\end{center}

\noindent \textbf{Due:} Friday, September 18 at 11:59pm on Canvas.\\ 

\noindent \textbf{Instructions:} There are two parts to this assignment. Part I is practice with maximum likelihood estimation, Part II is practice with logistic regression diagnostics.\\

\noindent \textbf{Getting started:} Begin by downloading the HW 3 template from the course website:\\

\url{https://sta214-f26.github.io/homework/hw_03_template.qmd}\\

\noindent Save this template file to your computer, then open it in RStudio. As you complete the assignment, you will write down your answers to all questions in the file, and include all R code in code chunks. \textit{If a question requires code, you will not receive credit if no code is provided.}\\

\noindent \textbf{Collaborators and AI statement:} At the top of your assignment, list any collaborators (Dr. Evans does not count -- he is a resource, not a collaborator), and provide an AI statement describing your use of AI to complete the questions. If you did not use AI in any way, let me know that too. If you did use AI, provide: the model used, the HW questions on which you used AI, and a summary of the ways in which AI provided assistance. As described in the syllabus, AI use is permitted on homework assignments, but submissions without a complete AI statement may not receive full credit.\\

\noindent \textbf{Submission:} When you have completed the assignment, render your homework to HTML and submit on Canvas. For Part I, you may submit a separate scan of your written work, if you prefer.\\

\noindent \textbf{Important:} you must submit the rendered HTML on Canvas to receive full credit!

\section{Maximum likelihood estimation}

This part focuses on practice with maximum likelihood estimation. A brief review of mathematical notation and rules for logs and derivatives can be found in the course codebook.\\

\noindent Suppose we have a random variable $Y$, and we assume that $Y \sim N(\mu, 1)$. (That is, $Y$ follows a normal distribution with  unknown mean $\mu$, and known variance $1$. The ``probability'' function for $Y$ is then
$$P(Y = y | \mu) = \frac{1}{\sqrt{2 \pi}} \exp\left\lbrace-\frac{1}{2}(y - \mu)^2\right\rbrace$$
(Here $\pi$ is the \textit{number}, i.e. $\pi = 3.14159...$). \\

\noindent (Note: Technically, this isn't really a probability in the same sense as probabilities for a Bernoulli distribution, but the distinction is not important for maximum likelihood estimation)\\

\noindent We get $n$ independent observations $Y_1,...,Y_n$ from this normal distribution, and we want to estimate $\mu$.

\begin{enumerate}
\item[1.] Show that the maximum likelihood estimator is $\widehat{\mu} = \frac{1}{n} \sum \limits_{i=1}^n Y_i$. Make sure to show all your work.
\end{enumerate}

\newpage

\section{Data analysis}

Here we work with data from a website called ScienceForums.Net (SFN), which has been open since 2002 and hosts conversations on a range of topics from biological and physical science to religion and philosophy. Each row in the data represents one ‘thread’, which is comprised of a series of posts stemming from an initial post. For each thread, we have some information that SFN collects such as the number of views and the number of authors. The threads present in the data are a random sample of threads from 2002-2014, with the data collected in 2014. SFN moderators are interested in using this data to determine which threads warrant the most attention.\\

\noindent You can load the SFN data into R by

\begin{verbatim}
sfn <- read.csv("https://sta214-f26.github.io/homework/sfn.csv")
\end{verbatim}

\noindent The \texttt{sfn} dataset contains the following columns:

\begin{itemize}
\item Age: the age of the thread (in days) when the data was collected in 2014, measured from the first post in the thread
\item State: sometimes moderators close threads if they are inappropriate. closed indicates the thread has been closed, otherwise State is open
\item Posts: the number of posts in the thread
\item Views: the total number of views of the thread
\item Duration: the number of days between the first and last posts in the thread
\item Authors: the number of distinct authors posting in the thread
\item AuthorExperience: the number of days the author of the first post in the thread had been registered on SFN when the thread began (0 indicates they registered that day)
\item DeletedPosts: the number of posts in the thread that have been deleted by a moderator
\item Forum: the forum in which the thread was posted (e.g., Science)
\item AuthorBanned: whether the original author of the thread is currently banned from posting on SFN (at the time of data collection, not when the thread was first posted)
\end{itemize}

\noindent \textbf{Research question:} Suppose you have been approached by moderators at SFN. They give you the data, and ask the following question:
\begin{itemize}
\item Is there a relationship between the number of Posts in a thread and whether a thread will have \textit{at least one} deleted post, after accounting for the number of Views, the number of Authors, and the Forum?
\end{itemize}

\begin{enumerate}
\item[2.] Here you will begin to use logistic regression to answer the moderators' question.

\begin{enumerate}
\item Which variables should we focus on to answer the moderators' question? Which of these is our response variable, and which will be our explanatory variables, for logistic regression? \textit{Hint: you may need to create a new variable for the response!}

\item Before fitting a model, let's see if any of the quantitative variables need transformations.
\begin{itemize}
\item Create empirical logit plots to summarize the relationship between each quantitative predictor and your binary response. See the class activity and the course codebook for details on empirical logit plots.
\item Using the empirical logit plots, discuss whether any transformations are needed on the explanatory variables.
\end{itemize}
\item Based on your empirical logit plots, write down a logistic regression model that will allow you to answer the moderators' question. Describe how you will use the model to answer their question.
\item Fit your model from (c), and report the equation of the fitted model.

\item Assess the shape assumption for your fitted model: create quantile residual plots to check the shape assumption for quantitative variables (you may use the \texttt{qresid} function in the \texttt{statmod} package). See the class activity and the course codebook for details on quantile residual plots.

\item Use the \texttt{regressinator} package (see the class activity and course codebook) to create a quantile residual plot lineup for one of your quantitative explanatory variables. Can you distinguish the original data from the simulated plots?

\item Based on your exploration in (e) and (f), do you think there are any violations to the shape assumption? If so, experiment with further transformations to address these violations, and report your final fitted model.  If there are no violations of the shape assumption, simply state this.

\item Interpret the coefficients of your final logistic regression model in context of the research question.

\end{enumerate} 
\end{enumerate}

\end{document}
